\documentclass[11pt,a4paper,twoside]{article}
\usepackage[utf8]{inputenc}
\usepackage[english]{babel}
\usepackage{actuarialangle, actuarialsymbol}
\usepackage{bbm}


\begin{document}

State the different forms of reinsurance and associated probabilities

$ X $ is the loss of the original insurance policy.
$ Y $ is the insurer’s retained loss after reinsurance.
$ Z $ is the reinsurer’s covered loss.

To find the associated probability distributions one can use the usual $ Y = g(x) \implies f_y(y) = f(g^{-1}(y))|\frac{dg^{-1}(y)}{dy}| $.
The expectation is usually easier to determine by LOTUS.

1. Quota-Share Reinsurance:
The insurer cedes a fixed proportion $ q $ of every loss to the reinsurer.
Insurer’s loss:
$
Y = (1 - q)X
$
Reinsurer’s loss:
$
Z = qX
$

2. Surplus Reinsurance:
The insurer retains losses up to a certain level $ r $, and the reinsurer covers the rest.
Insurer’s loss:
$ Y = \min(X, R) $
Reinsurer’s loss:
$Z = \max(0, X - R)$

3. Excess-of-Loss Reinsurance:
The reinsurer covers losses exceeding a retention $ r $, up to a limit $ m $.
Insurer’s loss:
$  Y = \min(X, a) + \max(0, X - (a + b)) $
Reinsurer’s loss:
$ Z = \max(0, \min(X - a, b)) $

The corresponding pdf's/cdf's are given by
1.
Insurer's loss:
PDF: $ f_Y(y) = \frac{1}{1 - q} f_X\left(\frac{y}{1 - q}\right) $
CDF: $ F_Y(y) = F_X\left(\frac{y}{1 - q}\right) $
Reinsurer’s loss:
PDF: $ f_Z(z) = \frac{1}{q} f_X\left(\frac{z}{q}\right) $
CDF: $ F_Z(z) = F_X\left(\frac{z}{q}\right) $
2.
Insurer's loss:
PDF:
\[
f_Y(y) =
\begin{cases}
f_X(y), & y \leq R, \\
0, & y > R
\end{cases}
\]
CDF:
\[
f_Y(y) =
\begin{cases}
F_X(y), & y \leq R, \\
1, & y > R
\end{cases}
\]
Reinsurer's loss
PDF:
$f_Z(z) = f_X(r + z), z \geq 0$
CDF:
$F_Z(z) = F_X(r + z), z \geq 0$
3.
Insurer's loss:
PDF:
\[
f_Y(y) =
\begin{cases}
f_X(y), & y < r, \\
F_X(r + m) - F_X(r), & y = r \\
f_X(y + m), & y > r \\
\end{cases}
\]
CDF:
\[
F_Y(y) =
\begin{cases}
F_X(y), & y < r, \\
F_X(m), & y = r \\
F_X(y + m), & y > r 
\end{cases}
\]
Reinsurer's loss
PDF:
\[
f_Z(y) =
\begin{cases}
F_X(r), & 0 = y, \\
f_X(y + r), & 0 < y \leq m, \\
0, & y > m
\end{cases}
\]
CDF:
\[
F_Z(y) =
\begin{cases}
F_X(r), & 0 = y, \\
F_X(y + r), & 0 < y < m, \\
1, & y \geq m
\end{cases}
\]


$f_Z(z) = f_X(r + z), z \geq 0$
CDF:
$F_Z(z) = F_X(r + z), z \geq 0$


State the formulas for probabilities and related objects in life insurance/survival analysis.

Complete future lifetime: \\
$T_x$ $\rightarrow$ The pdf of $T_x$ is $f_x(t) = \px[t]{x} \mu_{x+t}$. 
\newline
Curtate future lifetime: \\
$K_x = [T_x]$ $\rightarrow$ The pmf of $K_x$ is $\Pr(K_x = k) =  \px[k]{x} q_{x+k}$.
\newline
Survival function: $S_x(t) = \Pr(T_x > t)$ \\
Formulas: \\
$S_x(u+t) = S_x(u) S_{x+u}(t)$ $\rightarrow$ $S_{x+u}(t) = \frac{S_x(u+t)}{S_x(u)}$ \\
$S_0(x+t) = S_0(x) S(t)$ $\rightarrow$ $S(t) = \frac{S_0(x+t)}{S_0(x)}$
\newline

Three conditions: \\
(1) $S_x(0) = 1$ \\
(2) $\lim_{t \to \infty} S_x(t) = 0$ \\
(3) $S_x(t)$ must be a non-increasing function of $t$. \\
\newline
Acturial Notation: \\
Survival probability: $ \px[t]{x} = \Pr(T_x > t) $ \\
Mortality probability: $ \qx[t]{x} = 1 - \px[t]{x} = \Pr(T_x \leq t) $ \\
Formulas: \\
$ \px[t+u]{x} = \px[t]{x} \px[u]{x+t} $ \\
$ \qx[t+u]{x} = \qx[t]{x} + \px[t]t{x} \qx[u]{x+t}$ \\
$ \qx[t|u]{x} = \px[t]{x}\qx[u]{x+t} = \px[t]{x} - \px[t+u]{x} = \qx[t+u]{x} - \qx[t]{x} $ \\


State the formula's related to life tables

Number of lives: $ \lx{x} $ \\
Number of deaths: $ \dx[t]{x} = \lx{x} - \lx{x+t} $ \\
Formulas: \\
$ \px[t]{x} = \frac{\lx{x+t}}{\lx{x}} $ \\
$ \qx[t]{x} = \frac{\dx[t]{x}}{\lx{x}} = \frac{\lx{x} - \lx{x+t}}{\lx{x}} $ \\
$ \qx[t|u]{x} = \frac{\dx[u]{x+t}}{\lx{x}} = \frac{\lx{x + t} - \lx{x+t+u}}{\lx{x}} $

State the formula's related to the force of mortality

Definition: $ \mu_x(t) = \mu_{x+t} = \frac{f_x(t)}{S_x(t)} = - \frac{d}{dt}\log S_x(t) $ \\
where $ f_x(t) = \mathbb{P}(T(x) = t) $ and $ S_x $ the corresponding survival function.
Condition: $ \lim_{t \to \infty} \int_{0}^t \mu_s ds = \infty $ \\
Formulas: \\
$ \px[t]{x} = e^{-\int{0}^t \mu_{x+s}ds} = e^{-\int{x}^{x+t} \mu_s ds} $ \\
$ \qx[t]{x} = \int{0}^t \px[s]{x} \mu_{x+s}ds $ \\
$ \qx[t|u]{x} = \int_{t}^{t+u} \px[s]{x}\mu{x+s} ds $


Give expressions for the different kinds of expected future lifetime

Expectations: \\
Complete expectation of life (time (so not necessarily in units year) completed prior death)
$E[T_x] = \eringx{x} = \int_0^\infty t \px[t]{x} \mu_{x+t} dt = \int_0^\infty \px[t]{x} dt$ \\
Curtate Future Lifetime (number of years (so it's integer) completed prior death)
$E[K_x] = e_x = \sum_{k=1}^\infty k \px[k]{x} q_{x+k} = \sum_{k=1}^\infty \px[k]{x}$ \\
$E[\min(T_x, n)] = \eringx_{x:\angln} = \int_0^n t \px[t]{x} \mu_{x+t} dt + n \px[n]{x} = \int_0^n \px[t]{x} dt$ \\
$E[\min(K_x, n)] = e_{x:\angln} = \sum_{k=1}^{n-1} k \px[k]{x} q_{x+k} + n \px[n]{x} = \sum_{k=1}^n \px[k]{x}$ \\
\newline
Second moments: \\
$E[T_x^2] = \int_0^\infty t^2 \px[t]{x} \mu_{x+t} dt = \int_0^\infty 2t \px[t]{x} dt$ \\
$E[K_x^2] = \sum_{k=1}^\infty k^2 \px[k]{x} q_{x+k} = \sum_{k=1}^\infty (2k-1) \px[k]{x} = 2 \sum_{k=1}^\infty k \px[k]{x} - e_x$
\newline
Variance:
$Var(T_x) = E[T_x^2] - E[T_x]^2$ \\
$Var(K_x) = E[K_x^2] - E[K_x]^2$
\newline
Recursive formulas: \\
$\eringx{x}= \eringx{\endowxn} + \px[n]{x} \eringx{x+n}$ \\
$e_x = e_{x:\angln} + \px[n]{x} e_{x+n} \quad \rightarrow \quad e_x = p_x + p_x e_{x+1}$ \\
$\overset{\circ}{e}_{x:\overline{m}|} = \overset{\circ}{e}_{x+m:\overline{n-m}|} + m {}_m p_x \overset{\circ}{e}_{x+m:\overline{n-m}|} \text{ for } m < n$ \\
$e_{x:\overline{m}|} = e_{x+m:\overline{n-m}|} + m {}_m p_x e_{x+m:\overline{n-m}|} \text{ for } m < n \quad \rightarrow \quad e_{x:\overline{n}|} = p_x + p_x e_{x+1:\overline{n-1}|}$


State Wald's identity and derive it

Wald’s identity \\
Let $ T, X_1, X_2, \dots $ be independent real random variables in $ \mathcal{L}^1(\mathbb{P}) $. \\
Let $ \mathbb{P}[T \in \mathbb{N}_0] = 1 $ and assume that $ X_1, X_2, \dots $ are identically distributed. \\
Define
$$
S_T := \sum_{i=1}^{T} X_i.
$$
Then $ S_T \in \mathcal{L}^1(\mathbb{P}) $ and $ \mathbb{E}[S_T] = \mathbb{E}[T] \mathbb{E}[X_1] $. \\

\newline
Proof: Define $ S_n = \sum_{i=1}^{n} X_i $ for $ n \in \mathbb{N}_0 $. \\
Then 
$$
S_T = \sum_{n=1}^{\infty} S_n \mathbbm{1}_{(T=n)}.
$$
By Remark 2.15, the random variables $ S_n $ and $ \mathbbm{1}_{(T=n)} $ are independent for any $ n \in \mathbb{N} $ and thus uncorrelated. This implies (using the triangle inequality; see Theorem 5.3(v)) \\
$$
\mathbb{E}[S_T] = \sum_{n=1}^{\infty} \mathbb{E}[S_n \mathbbm{1}_{(T=n)}] 
= \sum_{n=1}^{\infty} \mathbb{E}[S_n] \mathbb{E}[\mathbbm{1}_{(T=n)}]
$$
$$
\leq \sum_{n=1}^{\infty} \mathbb{E}[X_1] n \mathbb{P}[T = n] = \mathbb{E}[X_1] \mathbb{E}[T].
$$
The same computation without absolute values yields the remaining part of the claim.

Derive Blackwell-Girshicks (Wald's) identity for compound variance 

Let $T, X_1, X_2, \dots$ be independent real random variables in $\mathcal{L}^2(\mathbb{P})$. \\
Let $\mathbb{P}[T \in \mathbb{N}_0] = 1$ and let $X_1, X_2, \dots$ be identically distributed. \\
Define 
$$ S_T := \sum_{i=1}^{T} X_i. $$  
Then $S_T \in \mathcal{L}^2(\mathbb{P})$ and  
$$ \textbf{Var}[S_T] = \mathbb{E}[X_1]^2 \text{Var}[T] + \mathbb{E}[T] \text{Var}[X_1]. $$
Proof: \\
Define $S_n = \sum_{i=1}^{n} X_i$ for $n \in \mathbb{N}$. Then (as in the proof of Wald’s identity) $S_n$ and $\mathbbm{1}_{(T=n)}$ are independent; \\
hence $S_n^2$ and $\mathbbm{1}_{(T=n)}$ are uncorrelated and thus \\
$$ \mathbb{E}[S_T^2] = \sum_{n=0}^{\infty} \mathbb{E}[\mathbbm{1}_{(T=n)} S_n^2] $$  
$$ = \sum_{n=0}^{\infty} \mathbb{E}[\mathbbm{1}_{(T=n)}] \mathbb{E}[S_n^2] $$  
$$ = \sum_{n=0}^{\infty} \mathbb{P}[T = n] (\text{Var}[S_n] + \mathbb{E}[S_n]^2) $$  
$$ = \sum_{n=0}^{\infty} \mathbb{P}[T = n] (n \text{Var}[X_1] + n^2 \mathbb{E}[X_1]^2) $$  
$$ = \mathbb{E}[T] \text{Var}[X_1] + \mathbb{E}[T^2] \mathbb{E}[X_1]^2. $$  

By Wald’s identity (Theorem 5.5), we have $\mathbb{E}[S_T] = \mathbb{E}[T] \mathbb{E}[X_1]$; hence 
$$ \text{Var}[S_T] = \mathbb{E}[S_T^2] - \mathbb{E}[S_T]^2 = \mathbb{E}[T] \text{Var}[X_1] + (\mathbb{E}[T^2] - \mathbb{E}[T]^2) \mathbb{E}[X_1]^2$$



\end{document}
